\clearpage\SourceLeaf{137}
\begin{center}
[III.]

\medskip
\textbf{WORKED CALCULATION. GIVEN $A=p\cos P$, $B=p\sin P$,\\
FIND $p$ AND $P$.}
\end{center}

\begin{center}\small
\begin{tabular}{l@{\;}l@{\qquad\hspace{.8em}}l}
$B\; .\ .\ .\ .\ .\ .$ & $9{,}56905\;69225\,n$ &\\
$A\; .\ .\ .\ .\ .\ .$ & $0{,}28523\;04177\,n$ &\\\cmidrule(r{1.4em}){1-2}
$\mathrm{tg}\,P\; .\ .\ .\ .\ .$ & $9{,}28382\;65048$ & $P=190^0\,52'\,52''{,}913567$\\
$\mathrm{tg}\,P^*\; .\ .\ .\ .$ & $9{,}28379\;34121$ & $P^*=190\;52\;50$\\\cmidrule(r{1.4em}){1-2}
$\cos P^*\; .\ .\ .$ & $9{,}99212\;15841\,n$ &\\
& $-11792$ & $5{,}5197322\; .\ .\ .\ .\ .\;3\;30927$\\
$\cos P\; .\ .\ .\ .$ & $9{,}99212\;04049\,n$ & $\underline{9{,}9842419}\; .\ .\ .\ .\;\cos P\cos P^*$\\\cmidrule(r{1.4em}){1-2}
$p\; .\ .\ .\ .\ .\ .\ .$ & $0{,}29311\;00128$ & $5{,}5039741$
\end{tabular}
\end{center}

\begin{center}\small
\renewcommand{\arraystretch}{1.35}
\begin{tabular}{r@{\;}l@{\;}l}
\multirow{2}{*}{$11792\Bigg\{$}&$9{,}2838100\; .\ .\ .\ .\;\BareRoot{(\mathrm{tg}\,P\mathrm{tg}\,P^*)}$&\\
&$\overline{4{,}7877841}$&\multirow{2}{*}{$\Bigg\}\;2''{,}913567.$}\\
$\left[-\log\frac{\mathrm{Mod.}}{206264{,}8}\right]=\mathrm{Const.}\; .\ .\ .$&$-4{,}3233592$&
\end{tabular}
\end{center}

\medskip
\begin{center}\textbf{REMARKS.}\end{center}

That the review on pp. 121--127 was written by \textsc{Gauss} is clear from a
letter from \textsc{Gauss} to \textsc{Olbers} dated 3 December 1808. The letter reads:

\begin{quotation}\itshape
„I should very much like, for example, a table of the kind proposed by \textup{\textsc{Leonelli}}
(my review: Hallische L.Z. 1808 of 12 Feb.) to be made. I once began one myself
to just 5 decimal places. In calculations that require very many logarithms of
sums or differences (as, for example, in my method for calculating perturbations),
such a table would considerably ease the work«.
\end{quotation}

Note [I.] is taken from a handbook. [II.] and [III.] are on sheets attached to the
books in the \textsc{Gauss} library by \textsc{Hobert} and \textsc{Ideler},
Trigonometrische Tafeln, and \textsc{Vega}, Thesaurus Logarithmorum.
\hfill \textsc{Börsch, Krüger}.
